% ============================================================
%  CHAPITRE II : État de l'art
% ============================================================
\chapter{État de l'art}
\label{chap:etat_art}

\section*{Introduction}

Ce chapitre présente l'analyse des référentiels normatifs mobilisés pour concevoir la méthodologie d'auto-évaluation de la plateforme. Contrairement à une étude comparative d'outils concurrents existants, ce chapitre est centré sur l'analyse du guide \gls{cmrpi}/\gls{ausim} et sa mise en correspondance avec les référentiels internationaux \gls{iso27001} et \gls{nistcsf}, ainsi que sur la justification académique des choix méthodologiques retenus pour le questionnaire.

% -------------------------------------------------------
\section{Référentiels normatifs analysés}
\label{sec:referentiels}

\subsection{Guide CMRPI/AUSIM}

\hspace{1cm}Le guide des bonnes pratiques \gls{cmrpi}/\gls{ausim}~\cite{cmrpiausim} est un document de référence d'une quarantaine de pages, élaboré conjointement par le \gls{cmrpi} et l'\gls{ausim}. Il couvre cinq domaines de bonnes pratiques en cybersécurité, alignés avec les standards internationaux présentés ci-après. C'est ce guide qui constitue la base de contenu du questionnaire d'auto-évaluation développé dans le cadre de ce stage.

\subsection{ISO/IEC 27001:2022}

L'\gls{iso27001}~\cite{iso27001} est le référentiel international de gestion de la sécurité de l'information le plus largement reconnu. Il définit les exigences relatives à la mise en place, à la mise en \oe{}uvre, à la maintenance et à l'amélioration continue d'un système de management de la sécurité de l'information. Dans le cadre de ce projet, les 93 contrôles de l'Annexe A de la norme, organisés en quatre thèmes (A.5 Organisationnel, A.6 Personnes, A.7 Physique, A.8 Technologique), ont été regroupés conceptuellement en cinq domaines afin de les rendre exploitables par un questionnaire court et compréhensible par un public non technique --- ce regroupement est détaillé en section~\ref{sec:domaines}.

\subsection{NIST Cybersecurity Framework 2.0}

Le \gls{nistcsf}~\cite{nistcsf2024} structure la gestion des risques cybersécurité autour de six fonctions : \textit{Govern} (gouverner), \textit{Identify} (identifier), \textit{Protect} (protéger), \textit{Detect} (détecter), \textit{Respond} (répondre) et \textit{Recover} (récupérer) -- la fonction \textit{Govern} ayant été ajoutée aux cinq fonctions d'origine (\textit{Identify}, \textit{Protect}, \textit{Detect}, \textit{Respond}, \textit{Recover}) lors de la publication de la version 2.0 du référentiel en février 2024. Une correspondance conceptuelle a été établie entre ces six fonctions et les cinq domaines du questionnaire, ce qui a permis de vérifier la cohérence et l'exhaustivité relative de la couverture retenue.

\subsection{Loi 09-08 (Maroc)}

La loi 09-08~\cite{loi0908} relative à la protection des personnes physiques à l'égard du traitement des données à caractère personnel encadre, au Maroc, la collecte et le traitement des données personnelles. Dans la plateforme développée, cette exigence réglementaire est intégrée sous la forme d'un consentement explicite recueilli via une case à cocher obligatoire lors de l'inscription de l'utilisateur, préalablement à toute collecte de données relatives à la \gls{pme} évaluée.

\iacsdefinition{Maturité cybersécurité}{
  La maturité cybersécurité désigne le degré auquel les pratiques, processus et contrôles de sécurité d'une organisation sont formalisés, appliqués de manière cohérente et intégrés à sa gouvernance. Elle se mesure généralement par paliers croissants, du niveau le plus faible (pratiques absentes ou ad hoc) au niveau le plus avancé (pratiques formalisées, mesurées et communiquées).
}

% -------------------------------------------------------
\section{Les cinq domaines retenus et leur justification}
\label{sec:domaines}

\subsection{Correspondance avec les référentiels internationaux}

Le tableau~\ref{tab:domaines} présente les cinq domaines retenus pour structurer le questionnaire d'auto-évaluation, ainsi que leur correspondance avec les sections de l'\gls{iso27001} et les fonctions du \gls{nistcsf}.

\begin{table}[H]
  \centering
  \caption{Les 5 domaines du questionnaire et leur correspondance normative}
  \label{tab:domaines}
  \renewcommand{\arraystretch}{1.4}
  \begin{tabularx}{\textwidth}{|c|X|c|c|}
    \hline
    \rowcolor{iacsTableHead}
    \textcolor{white}{\textbf{\#}} &
    \textcolor{white}{\textbf{Domaine}} &
    \textcolor{white}{\textbf{Réf.\ ISO 27001}} &
    \textcolor{white}{\textbf{Réf.\ NIST CSF}} \\
    \hline
    1 & Gouvernance & A.5.1--A.5.2 (Politiques, rôles) & Govern \\
    \hline
    \rowcolor{iacsLightGray}
    2 & Accès \& Identités & A.5.16--A.5.18, A.8.2, A.8.5 (Gestion des identités et des accès) & Protect \\
    \hline
    3 & Infrastructure \& Réseau & A.8.20--A.8.23 (Sécurité réseau) & Identify / Protect \\
    \hline
    \rowcolor{iacsLightGray}
    4 & Incidents \& Continuité & A.5.24--A.5.28, A.8.14 (Gestion des incidents, continuité) & Detect / Respond / Recover \\
    \hline
    5 & Sensibilisation \& Formation & A.6.3 (Sensibilisation et formation) & Protect \\
    \hline
  \end{tabularx}
\end{table}

\subsection{Justification académique du choix de 5 domaines et 25 questions}

\hspace{1cm}Le choix d'un questionnaire structuré en cinq domaines de cinq questions chacun, soit 25 questions au total, repose sur deux fondements académiques et une décision de conception assumée.

\textbf{Fondements issus du CMM/CMMI.} Le modèle de maturité par paliers proposé par Humphrey~\cite{humphrey1989} dans \textit{Managing the Software Process} définit cinq niveaux de maturité croissants. Cette structuration en cinq niveaux a directement inspiré la définition des cinq seuils de maturité utilisés dans le moteur de scoring de la plateforme (voir chapitre~\ref{chap:conception}, section~\ref{sec:scoring}).

\textbf{Échelle de Likert à cinq niveaux.} Les travaux de Dawes~\cite{dawes2008} sur la comparaison des échelles de Likert à 5, 7 et 10 points montrent qu'une échelle à cinq niveaux constitue un compromis optimal entre simplicité de réponse et granularité de l'information recueillie, en particulier pour un public non spécialiste. Ce résultat a motivé le choix d'une échelle de Likert à cinq niveaux pour chaque question du questionnaire, avec des libellés contextualisés propres à chaque question plutôt qu'une simple graduation numérique de 1 à 5.

\textbf{Décision de conception sur le nombre de questions.} Le choix de cinq questions par domaine, plutôt qu'un nombre plus élevé, résulte d'un arbitrage assumé entre la couverture de chaque domaine et la longueur totale du parcours pour un public non spécialiste (dirigeants de \gls{pme}) : au-delà du minimum généralement recommandé de trois items pour caractériser une dimension de manière stable, cinq questions par domaine permettent de couvrir les principales pratiques du guide \gls{cmrpi}/\gls{ausim} sans dépasser un total de 25 questions, jugé raisonnable pour un outil d'auto-évaluation complété en une seule session.

\iacsnote{Ces deux fondements académiques, complétés par la décision de conception exposée ci-dessus, justifient conjointement la structure retenue pour le questionnaire : 5 domaines $\times$ 5 questions, avec une échelle de réponse de Likert à 5 niveaux à libellés contextualisés. Cette structure conditionne directement la méthodologie de scoring présentée au chapitre~\ref{chap:conception}.}

\section*{Conclusion}

Ce chapitre a permis d'analyser les référentiels normatifs mobilisés pour la conception du questionnaire d'auto-évaluation --- le guide \gls{cmrpi}/\gls{ausim}, l'\gls{iso27001}, le \gls{nistcsf} et la loi 09-08 --- et de justifier académiquement la structure retenue en cinq domaines et 25 questions. Le chapitre suivant détaille la conception et la réalisation de la plateforme développée à partir de cette analyse.
